\documentclass[10pt,a4paper]{amsart}
\usepackage[margin=19mm]{geometry}
\usepackage{amsmath,amssymb,mathtools}
\usepackage[scr=rsfs,cal=euler]{mathalfa}
\usepackage{xfrac}
\usepackage[unicode,hidelinks,bookmarks=false]{hyperref}
\newcommand{\ie}{i.e.\ }
\newcommand{\cf}{cf.\ }
\newcommand{\resp}{respectively\ }
\newcommand{\Subsection}{\subsection}
\providecommand{\nobreakdash}{\nobreak\hbox{-}}
\setlength{\emergencystretch}{2em}
\multlinegap0pt
\usepackage{bbm}
\usepackage{enumerate}
\usepackage{amssymb}
\newtheorem{lem}{Lemma}
\title{Chebyshev's method applied to polynomials with two distinct roots}
\author{Tarakanta Nayak, Pooja Phogat}
\date{Source: arXiv:2512.20079v1 (CC BY 4.0)}
\begin{document}
\maketitle

\noindent\emph{Source and licence.} Chebyshev's method applied to polynomials with two distinct roots. Version arXiv:2512.20079v1 (CC BY 4.0), distributed by arXiv under the Creative Commons Attribution 4.0 International licence, http://creativecommons.org/licenses/by/4.0/, as recorded in the arXiv licence metadata for this exact version and shown on the abstract page https://arxiv.org/abs/2512.20079v1. Full licence notice: \texttt{licenses/arxiv-2512.20079-CC-BY-4.0.txt}.\par\smallskip

\noindent\textit{Editorial scope.} Counted unit: the article's lem criticalpointsC\_p (source0.tex lines 286--289). The article's complete original proof of it is delivered with it (source0.tex lines 290--311, one \texttt{proof} environment of 22 lines). No mathematics is written, repaired or re-derived: the statement and the proof are the article's own lines, in the article's own notation, and no retained line is substituted or re-flowed.  Retained uncounted as the source-local context the proof needs: lines 225--228.  Nothing was omitted from any retained span, so the package carries no omission record.  No retained line contains a citation, so the wrapper carries no bibliography and no bibliography entry.  No retained line contains a figure environment, an image-inclusion command or drawing code, so no image is delivered and the package declares no asset dependency.  Editorial formatting only, in addition: an AMS \texttt{amsart} preamble that restates the article's own declarations and macros used by the retained lines, namely the environments \texttt{lem} on separate counters, together with the standard packages those lines need.  The retained environments print in this document as Lemma 1 in order of appearance, whereas the article prints the same items with the numbering of its own full document.  The complete native source of the article as distributed in the arXiv e-print for this exact version is bundled unchanged as \texttt{sources/191534/source0.tex}.
\begin{equation}
	p(z) =z^k(z-1)^m ~\mbox{for}~k,m \geq 1.
\label{generalform-polynomial}	
\end{equation}  Then  $p'(z)=z^{k-1}(z-1)^{m-1}((k+m)z-k), p''(z)  =z^{k-2}(z-1)^{m-2}f(z)~\mbox{and}~ p'''(z) =z^{k-3}(z-1)^{m-3}g(z)$, where  $f(z)=(k+m)(k+m-1)z^2-2k(k+m-1)z+k(k-1) $ and $g(z)=(k+m)(k+m-1)(k+m-2)z^3-3k(k+m-1)(k+m-2)z^2+3k(k-1)(k+m-2)z-k(k-1)(k-2)$. 

\begin{lem}
There is no critical point of $C_p$ in $(-\infty, 0) \cup (1, \infty)$. Furthermore, $C_p '(x)>0$ for all $x \in (-\infty, 0) \cup (1, \infty)$. 
\label{criticalpointsC_p}
\end{lem}

\begin{proof}
Note that $$\begin{aligned}
	C_p^{\prime}(z) &=\frac{1}{2}\left(\frac{p(z)p''(z)}{\left(p'(z)\right)^2}\right)^2\left(3- \frac{p'(z)p'''(z)}{\left(p''(z)\right)^2}\right) =\left(p(z)\right)^2\left(\frac{\left(3\left(p''(z)\right)^2-p'(z)p'''(z)\right)}{2(p'(z))^4}\right)\\
	&=z^{2k}(z-1)^{2m}\left(\frac{3z^{2k-4}(z-1)^{2m-4}\left(f(z)\right)^2-z^{2k-4}(z-1)^{2m-4}\left((k+m)z-k \right)g(z)}{2(p'(z))^4}\right)\\
	&=\frac{3\left(f(z)\right)^2-\left(kz+mz-k \right)g(z)}{2\left(kz+mz-k \right)^4},
\end{aligned}
$$ where the polynomials $f,g $
 are as defined just after Equation~(\ref{generalform-polynomial}).
In other words,  $$C_p'(z)=\frac{F(z)}{2\left(kz+mz-k \right)^4},$$  where  
$F(z)=(k+m)^2 (k+m-1)(2k+2m-1)z^4 +4k(k+m)(k+m-1)(-2 k-2m+1)z^3+ 
3k(k(k+m-1)(3 k+3m -2)+(k-1)(k+m)^2)z^2 -2k(k-1)((4k+1)(k+m)-3k)z+k^2(k-1)(2k-1)$.
The critical points of  $C_p$ are $\frac{k}{k+m}$ (counted twice) and the four roots of the quartic polynomial $F$.

Note that the coefficients of the odd powers of $z$ in $F(z)$ are always negative, whereas those of the even powers are positive. Therefore, $F(z)$ is   positive for all $z \in (-\infty,0)$. Hence $C_p$ has no critical point in $(-\infty,0)$.
\par  In order to examine the behaviour of $F$ on  $(0,\infty)$, we consider
%$F(x+1)=(k+m)^2 (k+m-1)(2(k+m)-1)(x+1)^4 +4k(k+m)(k+m-1)(-2(k+m)+1)(x+1)^3+ 
%3k(k(k+m-1)(3(k+m)-2)+(k-1)(k+m)^2)(x+1)^2+k(k-1)(-2(4k+1)(k+m)+6k)(x+1)+k^2(k-1)(2k-1)$\\
%i.e.,\\
$F(x+1)=(k+m)^2(k+m-1)(2(k+m)-1)x^4+4m(k+m)(k+m-1)(2(k+m)-1)x^3+3m(k^2(4m-1)+km(8m-7)+2m(2m^2-3m+1))x^2+2m(m-1)(k(4m+1)+2m(2m-1))x+m^2(2m^2-3m+1)$.
For $k,m \geq 1$, each coefficient  of $F(x+1)$ is either zero or  positive. Moreover, the coefficient of $x^2$ is strictly greater than $1$. This gives that $F(x+1)>1$ for every $x>0$. Since every real number in $(1,\infty)$ can be written as $1+x$, for some $x>0$, it follows that $F((1,\infty))\subset (1,\infty)$. In other words,  $F$ has no zero in $(1,\infty)$, i.e., $C_p$ has no critical point in $(1,\infty)$.
\par It also follows from the two previous paragraphs that $C_p '(x)>0$ for all $x \in (-\infty, 0) \cup (1,  \infty)$.
\end{proof}

\end{document}
